% Faithful formal CCC2021 Section4 transcription; original notation and argument retained.
\section*{4 Perfect matching}
In this section we prove the following theorem:
\begin{authorclaim}{Theorem 2}
The size of any tree-like semantic $\mathrm{Res}(\oplus)$ refutation of the formula $\mathrm{PM}_{K_{n+2,n}}$ is $2^{\Omega(n)}$.
\end{authorclaim}
By Lemma 1, to prove Theorem 2 it is sufficient to show that $R^{1/3}_{\mathrm{pub}}(\oplus_2\operatorname{Search}(\mathrm{PM}_{K_{n+2,n}}))=\Omega(n)$.
Consider the communication problem $\oplus\mathrm{PM}_n^m$ that is defined as follows: Alice and Bob have matrices $X$ and $Y$ over $\mathbb F_2$ respectively, each of the matrices has size $m\times n$, where $m\ne n$. Their goal is to find an all-zero row or column or two 1-cells in the same row or column in the matrix $X+Y$.
\begin{authorclaim}{Proposition 15}
$R^{1/3}_{\mathrm{pub}}(\oplus_2\operatorname{Search}(\mathrm{PM}_{K_{n+2,n}}))\ge R^{1/3}_{\mathrm{pub}}(\oplus\mathrm{PM}_n^{n+2})$.
\end{authorclaim}
\begin{proof}
A Boolean matrix of size $(n+2)\times n$ naturally corresponds to a subset of edges of $K_{n+2,n}$. A falsified clause encoding that a vertex must be covered by a matching corresponds to an all-zero row or column of the matrix; a falsified clause, encoding that a vertex can not be covered by a matching twice, corresponds to two ones in the same row or column.
\end{proof}
Theorem 2 follows from Proposition 15 and the following theorem.
\begin{authorclaim}{Theorem 16}
$R^{1/3}_{\mathrm{pub}}(\oplus\mathrm{PM}_n^{n+2})=\Omega(n)$.
\end{authorclaim}
\begin{proof}
We assume that $n=4m+1$, where $m$ is a non-negative integer. If the theorem is true for all $n$ with the residue 1 modulo 4, then it also holds for all other $n$. Indeed, the protocol for $\oplus\mathrm{PM}_{n+1}^{n+3}$ can be used for $\oplus\mathrm{PM}_n^{n+2}$ by adding to Alice's matrix an extra column and a row with exactly one 1-cell on their intersection and to Bob's matrix an extra column and a row with all zeros.
Let $P_0$ be a protocol for $\oplus\mathrm{PM}_n^{n+2}$ transmitting at most $k$ bits. We are going to apply $P_0(X,Y)$ only to the instances where the matrix $X+Y$ does not contain all-zero rows or columns. Thus, we assume that with probability at least $2/3$ $P_0$ returns a tuple $(r_1,c_1,r_2,c_2)\in([n+2]\times[n])^2$ such that $(X+Y)_{r_1,c_1}=(X+Y)_{r_2,c_2}=1$ and either
\[
\begin{cases}r_1=r_2\\c_1\ne c_2\end{cases}
\quad\text{or}\quad
\begin{cases}r_1\ne r_2\\c_1=c_2.\end{cases}
\]
With $O(1)$ bits of communication Alice and Bob can verify whether the answer of $P_0$ is correct and return $\bot$ (failure) if it is not. Also, we can reduce the failure probability by the repetition of the protocol. Let $P$ be a protocol for $\oplus\mathrm{PM}_n^{n+2}$ under the promise that $X+Y$ does not contain all-zero rows and columns that uses $O(k)$ bits of communication and returns a correct answer with probability at least $99/100$ and $\bot$ otherwise.
We are going to construct a protocol for $\mathrm{DISJ}_m$ transmitting $O(k)$ bits, where $m=(n-1)/4$. Since by Theorem 8 any protocol for $\mathrm{DISJ}_m$ transmits $\Omega(m)$ bits, we conclude that $k=\Omega(m)$. Let Alice's input for $\mathrm{DISJ}_m$ be $a_1,\ldots,a_m$ and Bob's input be $b_1,\ldots,b_m$.
\begin{authorclaim}{Lemma 17}
There exist matrices $A(0),A(1),B(0),B(1)\in\mathbb F_2^{4\times4}$ such that $A(x)+B(y)$ is a permutation matrix iff $x\wedge y$ is 0 and
\begin{equation}
A(1)+B(1)=\begin{pmatrix}0&1&1&1\\1&0&0&0\\1&0&0&0\\1&0&0&0\end{pmatrix}.
\tag{1}
\end{equation}
\end{authorclaim}
\begin{proof}
We simply present matrices that satisfy the conditions:
\[
A(0)=\begin{pmatrix}0&1&1&0\\1&1&0&0\\1&0&1&0\\0&0&0&0\end{pmatrix};
\quad A(1)=\begin{pmatrix}0&0&0&0\\0&0&0&0\\0&0&0&0\\0&0&0&0\end{pmatrix};
\]
\[
B(0)=\begin{pmatrix}0&0&1&0\\0&1&0&0\\1&0&0&0\\0&0&0&1\end{pmatrix};
\quad B(1)=\begin{pmatrix}0&1&1&1\\1&0&0&0\\1&0&0&0\\1&0&0&0\end{pmatrix}.
\]
\end{proof}
Notice that Lemma 17 immediately allows to reduce $\mathrm{DISJ}_m$ to the problem of checking whether the sum of Alices and Bobs matrices is a permutation matrix. In order to achieve that, Alice builds a matrix $A=\operatorname{diag}(A(a_1),\ldots,A(a_m))$, Bob builds a matrix $B=\operatorname{diag}(B(b_1),\ldots,B(b_m))$. It is easy to see that $A+B$ is a permutation matrix iff $\mathrm{DISJ}_m(a,b)=1$.
Let us describe the reduction of $\mathrm{DISJ}_m$ to $\oplus\mathrm{PM}_n^{n+2}$. Alice and Bob first construct matrices $X_0$ and $Y_0$ of the following form:
\[
X_0=\begin{pmatrix}A&0_{(n-1)\times1}\\0_{1\times(n-1)}&1\\0_{1\times(n-1)}&1\\0_{1\times(n-1)}&1\end{pmatrix};
\quad Y_0=\begin{pmatrix}B&0_{(n-1)\times1}\\0_{1\times(n-1)}&0\\0_{1\times(n-1)}&0\\0_{1\times(n-1)}&0\end{pmatrix},
\]
then
\[
X_0+Y_0=\begin{pmatrix}A+B&0_{(n-1)\times1}\\0_{1\times(n-1)}&1\\0_{1\times(n-1)}&1\\0_{1\times(n-1)}&1\end{pmatrix},
\]
where $A+B$ is a permutation matrix iff $\mathrm{DISJ}_m(a,b)=1$. Then if $P(X_0,Y_0)$ returns two cells that do not belong to the column $n$ we may conclude that $\mathrm{DISJ}_m(a,b)=0$. If $P(X_0,Y_0)$ returns two cells from the $n$th column, then the value of $\mathrm{DISJ}_m(a,b)$ can not be uniquely determined. Notice that for $X_0$ and $Y_0$ constructed as above the protocol always returning $(n+1,n,n+2,n)$ solves $\oplus\mathrm{PM}_n^{n+2}$.
If $\mathrm{DISJ}_m(a,b)=0$, then the matrix $X_0+Y_0$ contains at least two columns with three ones and these columns are indistinguishable from each over. To make use of that, we randomly shuffle rows and columns.
We are going to construct a protocol $T$ for $\mathrm{DISJ}_m$ as follows: Alice and Bob choose permutations $\pi\in S_n$, $\tau\in S_{n+2}$ and a matrix $\Delta\in\mathbb F_2^{(n+2)\times n}$ uniformly at random. We define matrices $X_0^{\tau,\pi}$ and $Y_0^{\tau,\pi}$ from $\mathbb F_2^{(n+2)\times n}$ such that for each $i\in[n+2]$ and $j\in[n]$, $(X_0^{\tau,\pi})_{i,j}=(X_0)_{\tau(i),\pi(j)}$ and $(Y_0^{\tau,\pi})_{i,j}=(Y_0)_{\tau(i),\pi(j)}$. Alice and Bob run the protocol $P$ for inputs $X=X_0^{\tau,\pi}+\Delta$, $Y=Y_0^{\tau,\pi}+\Delta$. Notice that $X+Y=X_0^{\tau,\pi}+Y_0^{\tau,\pi}$, thus $X+Y$ can be obtained from $X_0+Y_0$ by shuffling rows and columns. If $P(X,Y)$ returns two cells from the column $\pi(n)$, Alice and Bob return 1, if $P(X,Y)$ returns two cells from other column or row, Alice and Bob return 0. If $P(X,Y)$ returns $\bot$, then Alice and Bob return $\bot$.
First notice that if $\mathrm{DISJ}_m(a,b)=1$, then $T$ returns a correct answer or $\bot$ with probability 1 (and the probability of $\bot$ is at most $1/100$), since in that case $X+Y$ has exactly one column with three 1-cells, each of the other columns and rows contains exactly one 1-cell. Let us fix $a,b\in\{0,1\}^m$ such that $\mathrm{DISJ}_m(a,b)=0$. We denote $p:=\Pr[T(a,b)=0]$, we will show that $p\ge99/200$. We can then increase this probability to $2/3$ by repeating the protocol twice (if $T(a,b)$ returns 0 at least once, we return 0, if $T(a,b)$ always return $\bot$, we return $\bot$, otherwise we return 1).
Let us describe random bits used by the constructed protocol $T$. First, we use random bits $r$ to run the protocol $P$. Second, we use random bits to generate $\pi,\tau$, and $\Delta$. Since $\mathrm{DISJ}_m(a,b)=0$, we can fix $i\in[m]$ such that $a_i=b_i=1$. In that case the submatrix of $X_0+Y_0$ formed by rows and columns with the indices $4(i-1)+1,4(i-1)+2,4(i-1)+3,4(i-1)+4$ coincides with the matrix (1). Let us denote by $\operatorname{col}(j)$ for $j\in[n]$ the set of all tuples $(x,j,y,j)\in([n+2]\times[n])^2$.
\begin{align*}
p&=\Pr_{\pi,\tau,\Delta,r}[P_r(X,Y)\notin\operatorname{col}(\pi(n))]
-\underbrace{\Pr_{\pi,\tau,\Delta,r}[P_r(X,Y)=\bot]}_{=:p_\bot}\\
&=1-\Pr_{\pi,\tau,\Delta,r}[P_r(X,Y)\in\operatorname{col}(\pi(n))]-p_\bot\\
&=1-\sum_{\pi_0,\tau_0}\Pr_{r,\Delta}[P_r(X_0^{\tau_0,\pi_0}+\Delta,Y_0^{\tau_0,\pi_0}+\Delta)\in\operatorname{col}(\pi_0(n))]
\Pr_{\pi,\tau}[\pi=\pi_0,\tau=\tau_0]-p_\bot .
\end{align*}
Observe that for fixed $\pi_0$ and $\tau_0$ the random variable $(X_0^{\tau_0,\pi_0}+\Delta,Y_0^{\tau_0,\pi_0}+\Delta)$ is uniformly distributed over the pairs of matrices with the sum $X_0^{\tau_0,\pi_0}+Y_0^{\tau_0,\pi_0}$. Let $\alpha\in S_n$ be the transposition swapping $n$ and $4(i-1)+1$. Let $\beta\in S_{n+2}$ be the permutation swapping $n$ and $4(i-1)+2$, $n+1$ and $4(i-1)+3$, $n+2$ and $4(i-1)+3$ (i.e. $\beta$ is a product of three transpositions). By the construction of $\alpha$ and $\beta$, $(X_0+Y_0)=(X_0^{\beta,\alpha}+Y_0^{\beta,\alpha})$, thus $(X_0^{\tau,\pi}+Y_0^{\tau,\pi})=(X_0^{\tau\circ\beta,\pi\circ\alpha}+Y_0^{\tau\circ\beta,\pi\circ\alpha})$ for every $\pi,\tau$. Thus the random variable $(X_0^{\tau_0\circ\beta,\pi_0\circ\alpha}+\Delta,Y_0^{\tau_0\circ\beta,\pi_0\circ\alpha}+\Delta)$ has the same distribution with $(X_0^{\tau_0,\pi_0}+\Delta,Y_0^{\tau_0,\pi_0}+\Delta)$, thus we can continue the sequence as follows:
\begin{align*}
p&=1-\sum_{\pi_0,\tau_0}\Pr_{r,\Delta}[P_r(X_0^{\tau_0\circ\beta,\pi_0\circ\alpha}+\Delta,Y_0^{\tau_0\circ\beta,\pi_0\circ\alpha}+\Delta)\in\operatorname{col}(\pi_0(n))]\Pr_{\pi,\tau}[\pi=\pi_0,\tau=\tau_0]-p_\bot\\
&=1-\sum_{\pi_0,\tau_0}\Pr_{r,\Delta}[P_r(X_0^{\tau_0,\pi_0}+\Delta,Y_0^{\tau_0,\pi_0}+\Delta)\in\operatorname{col}(\pi_0\circ\alpha^{-1}(n))]\Pr_{\pi,\tau}[\pi=\pi_0,\tau=\tau_0]-p_\bot\\
&=1-\Pr_{\pi,\tau,\Delta,r}[P_r(X,Y)\in\operatorname{col}((\pi\circ\alpha^{-1})(n))]-p_\bot\\
&\ge1-\Pr_{\pi,\tau,\Delta,r}[P_r(X,Y)\notin\operatorname{col}(\pi(n))]-p_\bot=1-p-p_\bot .
\end{align*}
Thus, $p\ge1-p-p_\bot$ and $p\ge(1-p_\bot)/2=99/200$.
\end{proof}
